\paragraph{Supplement to Finding A51: finite matrices and the missing rationality premise.}
Two elementary statements distinguish a finite matrix over a completed
group algebra from a fixed rational Frobenius determinant. The first is
an abstract module construction, with no asserted geometric realization.
The second gives a conditional sufficient premise for the exponent
criterion of A47. Neither supplies an actual curve tower or a pointed
primitive factor.

\emph{1. Completed coefficients, including semilinear Frobenius, permit
arbitrary large conductor exponents.}
Let $p\ne\ell$ be primes and put
\[
 \mathcal R=\varprojlim_n\mathbf Z_p[C_{\ell^n}],
 \qquad \gamma_n\text{ a compatible generator of }C_{\ell^n}.
 \tag{FX.1}
\]
The transition maps are induced by the cyclic-group quotients.
For $a\ge0$ define the component of $e_a\in\mathcal R$ at level $n$ by
\[
 e_a^{(n)}=
 \begin{cases}
 1,&n\le a,\\[1mm]
 \displaystyle\ell^{-(n-a)}
       \sum_{j=0}^{\ell^{n-a}-1}\gamma_n^{j\ell^a},&n>a.
 \end{cases}
 \tag{FX.2}
\]
These averages are compatible under the quotient maps. On a character
of exact order $\ell^b$, their values are $1$ for $b\le a$ and $0$
otherwise. Splitting the finite cyclic group over a finite unramified
extension of $\mathbf Z_p$ therefore proves
$e_ae_c=e_{\min(a,c)}$. Thus
\[
 \varepsilon_0=e_0,\qquad \varepsilon_a=e_a-e_{a-1}\ (a\ge1)
 \tag{FX.3}
\]
are orthogonal conductor idempotents. At level $n$ they vanish for
$a>n$ and sum to $1$ for $0\le a\le n$.

Choose any integers $b_a\ge2$. The expression
\[
 B=\sum_{a\ge0}p^{b_a}\varepsilon_a\in\mathcal R
 \tag{FX.4}
\]
is well defined, since it is a finite sum at each level. Multiplication
by $1-B$ is an automorphism of the free rank-one $\mathcal R$-module.
For every exact-conductor-$\ell^a$ character its value is
$1-p^{b_a}$, a $p$-adic unit. The corresponding cokernel of
$1-(1-B)$ has exact exponent $p^{b_a}$.

This remains an obstruction to a bound when the deck action is
nonconstant. Fix $u\in\mathbf Z_\ell^\times$, acting on $\mathcal R$
by $\gamma\mapsto\gamma^u$. It fixes each $\varepsilon_a$.
The automorphism
\[
           \mathcal F=(1-B)u
 \tag{FX.5}
\]
is semilinear over $\mathcal R$ and linear over $\mathbf Z_p$.
After a finite unramified coefficient extension splitting the
level-$a$ cyclic group, its exact-conductor characters form $u$-orbits
of length $d_a=\operatorname{ord}_{\ell^a}(u)$ for $a\ge1$
(and $d_0=1$). On each orbit
$\mathcal F$ is a cyclic permutation times $c_a=1-p^{b_a}$.
The relations for the additive cokernel of $1-\mathcal F$ eliminate
all but one basis vector and leave the relation
$1-c_a^{d_a}$. Consequently the exact-conductor additive cokernel
has exponent
\[
          p^{b_a+v_p(d_a)}.
 \tag{FX.6}
\]
Here is the valuation check, including $p=2$. If $v_p(x)\ge2$,
then a power prime to $p$ preserves the valuation of $(1+x)^m-1$:
the linear term has valuation $v_p(x)$ and all later terms have
larger valuation. Raising $1+x$ to the $p$th power instead increases
that valuation by one, because its linear term has valuation
$v_p(x)+1$ and every other term has larger valuation. Apply these
two observations to $x=-p^{b_a}$ and factor $d_a$ into its $p$-part
and prime-to-$p$ part. Finite unramified scalar extension preserves
the elementary-divisor exponents, proving (FX.6) over $\mathbf Z_p$
as well. Since the $b_a$ are arbitrary, finite matrix size and this
semilinear structure alone allow arbitrarily rapid exponent growth.
No Weil-weight, polarization, Jacobian, or point-section condition
is claimed for these modules. This does not contradict the actual
descent identities of A45 or the conditional bounds in A51.

\emph{2. A fixed nonzero rational function has eventually constant
cross-prime valuation.}
Let $K/\mathbf Q_p$ be finite, with residue field
$\mathbf F_{p^s}$, uniformizer $\varpi$, and valuation
$v_\varpi(\varpi)=1$. For nonzero $f(X)\in K(X)$ write
\[
             f(X)=\varpi^\mu G(X)/H(X),
 \tag{FX.7}
\]
where $G,H$ are integral Laurent polynomials, each with a unit
coefficient. Clearing monomials does not change the valuations
at roots of unity; let $D$ bound the degrees of their nonzero
reductions after doing so. If $\zeta$ has exact order $\ell^a$ and
\[
                  \operatorname{ord}_{\ell^a}(p^s)>D,
 \tag{FX.8}
\]
then
\[
                        v_\varpi(f(\zeta))=\mu.
 \tag{FX.9}
\]
Indeed adjoining prime-to-$p$ roots of unity is unramified, and
reduction preserves their exact orders. The minimal polynomial of
$\bar\zeta$ over $\mathbf F_{p^s}$ has degree
$\operatorname{ord}_{\ell^a}(p^s)$. It cannot divide either nonzero
reduction of degree at most $D$, so both evaluated polynomials are
units. Finally $\operatorname{ord}_{\ell^a}(p^s)\to\infty$:
otherwise arbitrarily large $\ell^a$ would divide one of the
finitely many nonzero integers $p^{sj}-1$ with $j$ bounded.
This proves (FX.9) for every sufficiently large $a$ and every
primitive root of that level.

In particular, suppose an actual fixed integral Laurent matrix
$M(X)$ has nonzero determinant $f(X)=\det(1-M(X))$.
At high conductor the sum of the nonnegative elementary-divisor
valuations of $1-M(\zeta)$ is the constant $\mu$ in (FX.9).
Its cokernel therefore has uniformly bounded $p$-primary exponent.
Only an additional identification of the actual rational primitive
groups with such evaluated cokernels, or subgroups of finite direct
sums of them, would give a uniform $M_n$ for A47. In genus two and
characteristic two, along odd prime powers $n=\ell^a$, a uniform
$M_n$ would eventually satisfy $n>(2M_n-1)^2$.
No such geometric identification is obtained here.

\emph{3. Exact primary scope and the unresolved premise.}
Finite presentation of the full proper etale fundamental group is
known \cite[Theorem~1.1]{ESS2022}
\cite[Theorem~1.1 and Corollary~1.2]{Shusterman2022}, but the
presentations are profinite, not asserted to have discrete finite-word
relators or finite-word arithmetic Frobenius images. The familiar
surface-word presentation is for the prime-to-$p$ quotient
\cite[p.~4]{ESS2022}; its open subgroups have zero continuous
$H^1(-,\mathbf F_p)$ and do not retain the full curve's
characteristic-$p$ etale cohomology. Arithmetic Fox constructions in
the inspected number-field pro-$\ell$ setting have completed
power-series coefficients, including explicit infinite Milnor sums
\cite[Sections~4.1--4.3, Proposition~4.3.4]{KMT2017}.

The geometric skew-abelian tower paper of Asvin~G. states a finite
matrix description over a $\mathbf Z_\ell$ power-series ring.
Its return matrix is the product
around the whole growing Frobenius character orbit
\cite[Definition~8, Lemma~17, Section~3.1, equation~(6)]{Asvin2023}.
The congruences are at $\ell$, and its polynomial truncation uses
$\zeta_{\ell^a}-1$ being $\ell$-adically small
\cite[Theorems~23,26 and Section~3.3, p.~2169]{Asvin2023}.
That element is a $p$-adic unit for $p\ne\ell$. Neither this
truncation nor a determinant of a growing orbit product supplies
the single fixed rational function required in (FX.7).
There is also an integral issue in its geometric argument:
Lemma~11 claims that the natural pullback on integral $H^1$
is a direct factor for every finite curve map, and its proof
infers surjectivity on Tate modules from a surjective map of
abelian varieties. For
$[\ell]:E\to E$ on an elliptic curve, both induced maps are
multiplication by $\ell$ and have cokernel
$(\mathbf Z/\ell\mathbf Z)^2$. Thus that inference, and the
claimed natural-image saturation, fail integrally. We do not
use its geometric integral finite-free premise as an established
input. The abstract Section~3 coefficient-prime scope and the
two independent algebra proofs above do not depend on that premise.
Sinnott's rational-measure theorem likewise assumes an actual fixed
rational Fourier function and an appropriate unit-group support
condition \cite[Section~1.16 and Theorem~3.2]{Sinnott1987}.
No identification with that transform is provided by finite
profinite presentation or matrix size. These are completed
algebraic checks on the proposed route, not a geometric
counterexample or a closure of the primitive-factor supply gate.
